\newpage
\section{问题分析}

\subsection{问题一分析}
本题要求分析车主停车行为并建立选择概率模型。车主在选择停车场时，主要考虑因素包括：
\begin{enumerate}
    \item[(1)] \textbf{停车费用}：直接成本，影响用户决策的关键因素
    \item[(2)] \textbf{距离因素}：步行距离或等待时间
    \item[(3)] \textbf{便利性}：停车场位置、进出便捷程度
\end{enumerate}
采用MNL（Multinomial Logit）模型刻画用户选择行为。设，第：
U_i = \alpha - \beta c_i + \epsilon_i
其中，$\epsilon_i。选择概率为：
P_i = \frac{e^{U_i}}{\sum_{j=1}^{J} e^{U_j}} = \frac{e^{\alpha - \beta c_i}}{\sum_{j=1}^{J} e^{\alpha - \beta c_j}}

\subsection{问题二分析}
本题需要建立Stackelberg主从博弈模型。这是一个分层决策问题：
\begin{itemize}
    \item \textbf{领导者}：停车场运营商，制定价格策略最大化收益
    \item \textbf{追随者}：车主，根据价格选择最优点火场
\end{itemize}
Stackelberg均衡要求领导者的策略对追随者的最优反应是最优的。运营商收益为：
\max_{p} \pi_1(p) = \sum_{i=1}^{J} (p_i - c_i^{cost}) \cdot D_i(p)
其中(p) = D_{total} \cdot P_i(p)。

\subsection{问题三分析}
灵敏度分析是验证模型稳健性的重要手段。我们需要分析关键参数（如价格弹性系数$\beta$）变化对最优定价的影响。定义灵敏度系数：
S_{\beta}^{p^*} = \frac{\beta}{p^*} \cdot \frac{\partial p^*}{\partial \beta}
